\section{Первообразная}

\begin{definition}[первообразная]
Функция $F(x)$ называется \emph{первообразной} для функции $f(x)$ на интервале $(a,b)$, если
\[
  \forall x\in(a,b):\qquad F'(x)=f(x).
\]
\end{definition}

\begin{example}
Для $f(x)=\cos x$ первообразными служат $F_1(x)=\sin x$, $\;F_2(x)=\sin x+42$ и т.д.
\end{example}

\begin{remark}\leavevmode
\begin{enumerate}[label=\arabic*),leftmargin=*]
  \item Первообразная $F(x)$ дифференцируема, а значит \emph{непрерывна} на $(a,b)$.
  \item Первообразная задаётся именно на \emph{промежутке}: вне связного промежутка добавляемая
        константа на разных кусках может быть разной.
\end{enumerate}
\end{remark}

\begin{theorem}[об общем виде первообразных]
Если $F(x)$ --- какая-либо первообразная для $f(x)$ на $(a,b)$, то множество всех первообразных
для $f$ на этом интервале есть
\[
  \bigl\{\,F(x)+C \;:\; C\in\R\,\bigr\}.
\]
\end{theorem}

\begin{proof}\leavevmode
\begin{enumerate}[label=\arabic*),leftmargin=*]
  \item Любая функция вида $F(x)+C$ — первообразная: $\bigl(F(x)+C\bigr)'=F'(x)=f(x)$.
  \item Пусть $\Phi(x)$ — произвольная другая первообразная. Рассмотрим
        $\psi(x)=\Phi(x)-F(x)$. Тогда
        \[
          \psi'(x)=\Phi'(x)-F'(x)=f(x)-f(x)\equiv 0 \quad\text{на }(a,b).
        \]
        По следствию из теоремы Лагранжа функция с нулевой производной на промежутке постоянна,
        поэтому $\psi(x)\equiv C$, то есть $\Phi(x)=F(x)+C$. \qedhere
\end{enumerate}
\end{proof}

\paragraph{Геометрическая интерпретация.}
График любой первообразной получается из графика $F$ вертикальным сдвигом на константу $C$ ---
получается \emph{однопараметрическое семейство} «параллельных» кривых, целиком заполняющее полосу
$\{(x,y):x\in(a,b)\}$. Через каждую точку этой полосы проходит ровно одна кривая семейства;
касательные ко всем кривым в точках с одной и той же абсциссой $x$ имеют одинаковый угловой
коэффициент $f(x)$.

\section{Неопределённый интеграл}

\begin{definition}[неопределённый интеграл]
\emph{Неопределённым интегралом} функции $f(x)$ на $(a,b)$ называется совокупность всех её
первообразных. Обозначение:
\[
  \int f(x)\dx = F(x)+C,
\]
где $F$ — какая-либо первообразная, а $C$ пробегает $\R$.
\end{definition}

\subsection{Таблица основных интегралов}
\begin{align*}
&\textbf{Степенные и логарифм:} &
   \int x^{n}\dx &=\frac{x^{n+1}}{n+1}+C\ (n\neq-1), &
   \int\frac{dx}{x}&=\ln\abs{x}+C.\\[2pt]
&\textbf{Показательные:} &
   \int a^{x}\dx&=\frac{a^{x}}{\ln a}+C\ (a>0,\,a\neq1), &
   \int e^{x}\dx&=e^{x}+C.\\[2pt]
&\textbf{Тригонометрические:} &
   \int\sin x\dx&=-\cos x+C, & \int\cos x\dx&=\sin x+C,\\
&& \int\frac{dx}{\cos^2 x}&=\tg x+C, & \int\frac{dx}{\sin^2 x}&=-\ctg x+C.\\[2pt]
&\textbf{Гиперболические:} &
   \int\sh x\dx&=\ch x+C, & \int\ch x\dx&=\sh x+C.\\[2pt]
&\textbf{«Обратные»:} &
   \int\frac{dx}{a^2+x^2}&=\tfrac1a\arctg\tfrac xa+C, &
   \int\frac{dx}{\sqrt{a^2-x^2}}&=\arcsin\tfrac xa+C.
\end{align*}
\begin{align*}
  \int\frac{dx}{x^2-a^2}&=\frac1{2a}\ln\abs[\Big]{\frac{x-a}{x+a}}+C
     &&\text{(«высокий» логарифм)},\\[4pt]
  \int\frac{dx}{\sqrt{x^2\pm k}}&=\ln\abs[\big]{x+\sqrt{x^2\pm k}}+C
     &&\text{(«длинный» логарифм)}.
\end{align*}

\subsection{Линейность неопределённого интеграла}
\begin{property}[линейность]
Пусть $f,g$ имеют первообразные на $(a,b)$ и $\alpha,\beta\in\R$. Тогда
\[
  \int\bigl(\alpha f(x)+\beta g(x)\bigr)\dx
   =\alpha\int f(x)\dx+\beta\int g(x)\dx .
\]
\end{property}

\begin{proof}
Обозначим разность двух частей через
\[
  H(x)=\int\bigl(\alpha f+\beta g\bigr)\dx-\Bigl(\alpha\int f\dx+\beta\int g\dx\Bigr).
\]
Дифференцируя и пользуясь равенством $\bigl(\int\varphi\dx\bigr)'=\varphi$, получаем
\[
  H'(x)=\bigl(\alpha f+\beta g\bigr)-\bigl(\alpha f+\beta g\bigr)=0
  \quad\text{на }(a,b).
\]
Значит $H(x)\equiv\Const$, то есть левая и правая части отличаются лишь на константу — а с
точностью до константы неопределённый интеграл и определён. Следовательно, как семейства
первообразных они совпадают.
\end{proof}

\section{Взаимосвязь дифференцирования и интегрирования}

Операции дифференцирования и интегрирования взаимно обратны (с точностью до константы):
\begin{align*}
  \frac{d}{dx}\!\left(\int f(x)\dx\right) &= f(x), &
  d\!\left(\int f(x)\dx\right) &= f(x)\dx, &
  \int dF(x) &= F(x)+C.
\end{align*}
Первое равенство выражает, что интеграл «снимается» производной; третье — что интегрирование
дифференциала восстанавливает функцию с точностью до константы.

\section{Замена переменной}

В основе метода лежит формула производной сложной функции:
$\bigl[F(g(x))\bigr]'=F'(g(x))\,g'(x)$.

\begin{theorem}[о замене переменной]
Пусть $x=\varphi(t)$ определена и дифференцируема на интервале $T$, а её множество значений
$\varphi(T)$ лежит в области, где функция $f$ имеет первообразную $F$. Тогда
\[
  \int f(x)\dx\,\Big|_{x=\varphi(t)} \;=\; \int f\bigl(\varphi(t)\bigr)\varphi'(t)\dt,
  \qquad\text{то есть}\qquad
  F\bigl(\varphi(t)\bigr)+C=\int f\bigl(\varphi(t)\bigr)\varphi'(t)\dt .
\]
\end{theorem}

\begin{proof}
Достаточно проверить, что $F(\varphi(t))$ — первообразная для подынтегральной функции правой
части. Дифференцируя по $t$ по правилу дифференцирования сложной функции, получаем
\[
  \frac{d}{dt}\,F\bigl(\varphi(t)\bigr)=F'\bigl(\varphi(t)\bigr)\varphi'(t)
   =f\bigl(\varphi(t)\bigr)\varphi'(t).
\]
Итак, $F(\varphi(t))$ — первообразная для $f(\varphi(t))\varphi'(t)$, откуда и следует доказываемая
формула.
\end{proof}

\begin{example}
$\displaystyle\int\cos(5x+4)\dx$. Полагаем $t=5x+4$, тогда $dt=5\dx$, $\dx=\tfrac15\dt$:
\[
  \int\cos(5x+4)\dx=\int\cos t\cdot\tfrac15\dt=\tfrac15\sin t+C=\tfrac15\sin(5x+4)+C.
\]
\end{example}

\begin{example}
$\displaystyle\int\frac{x\dx}{1+x^2}$. Заметим $x\dx=\tfrac12\,d(x^2+1)$:
\[
  \int\frac{x\dx}{1+x^2}=\frac12\int\frac{d(x^2+1)}{x^2+1}=\frac12\ln(x^2+1)+C.
\]
\end{example}

\section{Интегрирование по частям}

Исходим из формулы дифференциала произведения $d(uv)=u\,dv+v\,du$.

\begin{theorem}[интегрирование по частям]
Если $u(x)$ и $v(x)$ дифференцируемы на интервале, то
\[
  \int u\,dv = uv-\int v\,du,
  \qquad\text{или}\qquad
  \int u(x)v'(x)\dx=u(x)v(x)-\int u'(x)v(x)\dx.
\]
\end{theorem}

\begin{proof}
Проинтегрируем тождество $d(uv)=u\,dv+v\,du$:
$\int d(uv)=\int u\,dv+\int v\,du$, то есть $uv=\int u\,dv+\int v\,du$, откуда $\int u\,dv=uv-\int v\,du$.
\end{proof}

\paragraph{Основные типы интегралов, берущихся по частям.}
\begin{enumerate}[label=\textbf{(\alph*)},leftmargin=*]
  \item $\displaystyle\int P(x)e^{ax}\dx,\ \int P(x)\sin bx\dx,\ \int P(x)\cos bx\dx$ — берём $u=P(x)$
        (понижаем степень многочлена);
  \item $\displaystyle\int P(x)\ln x\dx,\ \int P(x)\arctg x\dx,\ \int P(x)\arcsin x\dx$ — берём $u$
        равным $\ln x$, $\arctg x$ и т.п. (избавляемся от «неудобной» функции);
  \item $\displaystyle\int e^{ax}\sin bx\dx,\ \int e^{ax}\cos bx\dx$ — \emph{возвратные}: после двух
        интегрирований по частям возникает исходный интеграл.
\end{enumerate}

\begin{example}
$\displaystyle\int x\,e^{x}\dx$: берём $u=x$, $dv=e^x\dx$, тогда $du=\dx$, $v=e^x$:
\[
  \int x e^{x}\dx=xe^{x}-\int e^{x}\dx=xe^{x}-e^{x}+C.
\]
\end{example}

\begin{example}
$\displaystyle\int\ln x\dx$: $u=\ln x$, $dv=\dx$, $du=\tfrac{dx}{x}$, $v=x$:
\[
  \int\ln x\dx=x\ln x-\int x\cdot\tfrac1x\dx=x\ln x-x+C.
\]
\end{example}

\begin{example}[возвратный интеграл]
$\displaystyle I=\int e^{x}\cos x\dx$. Дважды интегрируем по частям:
\[
  I=e^{x}\sin x-\int e^{x}\sin x\dx
   =e^{x}\sin x-\Bigl(-e^{x}\cos x+\int e^{x}\cos x\dx\Bigr)
   =e^{x}(\sin x+\cos x)-I.
\]
Отсюда $2I=e^{x}(\sin x+\cos x)$, то есть $I=\tfrac12 e^{x}(\sin x+\cos x)+C$.
\end{example}

\section{Комплексный метод в вещественном анализе}

Тот же интеграл удобно считать через комплексную экспоненту, опираясь на формулу Эйлера
$e^{i\varphi}=\cos\varphi+i\sin\varphi$. Тогда
\[
  e^{ax}\cos bx=\Re\,e^{(a+ib)x},\qquad e^{ax}\sin bx=\Im\,e^{(a+ib)x}.
\]
Интеграл от комплексной экспоненты берётся элементарно:
\[
  \int e^{(a+ib)x}\dx=\frac{1}{a+ib}\,e^{(a+ib)x}+C
   =\frac{a-ib}{a^2+b^2}\,e^{ax}\bigl(\cos bx+i\sin bx\bigr)+C.
\]
Выделяя вещественную и мнимую части, получаем сразу обе формулы:
\begin{align*}
  \int e^{ax}\cos bx\dx &= \Re\!\int e^{(a+ib)x}\dx
     =\frac{e^{ax}\bigl(a\cos bx+b\sin bx\bigr)}{a^2+b^2}+C,\\[4pt]
  \int e^{ax}\sin bx\dx &= \Im\!\int e^{(a+ib)x}\dx
     =\frac{e^{ax}\bigl(a\sin bx-b\cos bx\bigr)}{a^2+b^2}+C.
\end{align*}
При $a=b=1$ отсюда снова $\int e^{x}\cos x\dx=\tfrac12 e^{x}(\cos x+\sin x)+C$ — согласуется с
«возвратным» вычислением выше.

\section{Дифференцирование интеграла по параметру}

\begin{theorem}[правило Лейбница]
Если $f(x,\alpha)$ и $f'_{\alpha}(x,\alpha)$ непрерывны на прямоугольнике
$[a,b]\times[\alpha_1,\alpha_2]$, то функция $\Phi(\alpha)=\displaystyle\int_a^b f(x,\alpha)\dx$
дифференцируема по $\alpha$ и
\[
  \Phi'(\alpha)=\frac{d}{d\alpha}\int_a^b f(x,\alpha)\dx=\int_a^b f'_{\alpha}(x,\alpha)\dx .
\]
\end{theorem}

\noindent Для неопределённого интеграла это даёт удобный приём: продифференцировав известную формулу
по параметру, получаем новую.

\begin{example}[простейшее применение]
Из $\displaystyle\int e^{\alpha x}\dx=\tfrac1\alpha e^{\alpha x}+C$ дифференцированием по $\alpha$:
\[
  \frac{d}{d\alpha}\!\int e^{\alpha x}\dx=\int x\,e^{\alpha x}\dx,
  \qquad
  \frac{d}{d\alpha}\Bigl(\tfrac1\alpha e^{\alpha x}\Bigr)=-\tfrac1{\alpha^2}e^{\alpha x}+\tfrac{x}{\alpha}e^{\alpha x},
\]
откуда $\displaystyle\int x\,e^{\alpha x}\dx=-\tfrac1{\alpha^2}e^{\alpha x}+\tfrac{x}{\alpha}e^{\alpha x}+C$.
\end{example}

\subsection{Интегралы вида $\displaystyle\int\frac{dx}{(x^2+a^2)^n}$}
Отправная формула —
\[
  I_1(a)=\int\frac{dx}{x^2+a^2}=\frac1a\arctg\frac xa+C .
\]
Дифференцируем её по параметру $a$. Слева, поскольку
$\dfrac{\partial}{\partial a}\dfrac{1}{x^2+a^2}=-\dfrac{2a}{(x^2+a^2)^2}$, получаем
\[
  \frac{d}{da}\,I_1(a)=-2a\int\frac{dx}{(x^2+a^2)^2}.
\]
Справа
\[
  \frac{d}{da}\!\left(\frac1a\arctg\frac xa\right)
   =-\frac1{a^2}\arctg\frac xa+\frac1a\cdot\frac{1}{1+(x/a)^2}\cdot\Bigl(-\frac{x}{a^2}\Bigr)
   =-\frac1{a^2}\arctg\frac xa-\frac{x}{a\,(x^2+a^2)} .
\]
Приравнивая и выражая искомый интеграл, получаем рекуррентный результат:
\[
  \boxed{\;\int\frac{dx}{(x^2+a^2)^2}
   =\frac{1}{2a^3}\arctg\frac xa+\frac{x}{2a^2\,(x^2+a^2)}+C\;}
\]
Повторное дифференцирование по $a$ позволяет последовательно понижать степень и получать
$\displaystyle\int\frac{dx}{(x^2+a^2)^n}$ при любом $n$.
